\chapter{Estado del Arte}

Reconstruir imágenes a partir de la actividad cerebral no es un tema nuevo; lo que sí ha 
cambiado, y bastante, es cómo se aborda. Al principio se buscaba recuperar los píxeles de 
manera directa. Hoy el mismo problema se entiende de otra forma: como una traducción semántica 
entre dos mundos, el del córtex y el de los espacios latentes multimodales. Este capítulo 
sigue esa evolución, desde las primeras redes convolucionales hasta lo que se hace ahora con 
difusión y aprendizaje contrastivo.

\section{Primera Generación: enfoques convolucionales y bayesianos}

Los primeros trabajos que sacaron resultados decentes con fMRI recurrieron a redes 
convolucionales preentrenadas, usándolas como extractores de características.

\subsection{Regresión a espacios VGG}

Shen et al.~\cite{shen2019deep} y Beliy et al.~\cite{beliy2019voxels} encontraron que la 
actividad de V1–V3 tiene buena correlación con las primeras capas de VGG-19 y AlexNet. 
La receta era casi siempre la misma: armas una regresión lineal (Ridge o dispersa) desde los
vóxeles hacia las activaciones de la red, y después das la vuelta a la red con un
Deep Image Prior para construir una imagen que pegue con esas activaciones.

\subsection{Decodificación bayesiana y espacios discretos (VQGAN)}

Koide-Majima et al.~\cite{koidemajima2024} llevaron la misma línea más allá con un marco
bayesiano sobre el espacio latente discreto de VQGAN. Mezclando un componente perceptual 
con uno semántico lograron una identificación pareada bastante buena.
El problema de fondo de esta primera generación es de cómo está armada:
VGG es un espacio puramente visual y VQGAN es un diccionario discreto, y
ninguno conoce de semántica. Por eso las reconstrucciones quedaban parecidas
en la estructura general pero con significados aleatorios: una pelota con textura
de perro, una calle con objetos sin nombre. Por eso en esta tesis VQGAN+VGG
aparece sólo como referencia histórica, no es la arquitectura que se corre.

\section{Segunda Generación: difusión latente}

Cuando salió Stable Diffusion la cosa cambió en pocos meses.

\subsection{Condicionamiento dual}

Takagi y Nishimoto~\cite{takagi2023} fueron los primeros en usar Stable Diffusion
como decodificador directo. En lugar de predecir activaciones convolucionales
gigantes, entrenaron mapeos lineales hacia dos blancos más chicos: la
representación latente del VAE y el embedding de texto del estímulo.

\subsection{Brain-Diffuser}

Ozcelik y VanRullen~\cite{ozcelik2023} llevaron la idea a una difusión versátil condicionada con
embeddings visuales continuos predichos desde fMRI. El salto en calidad se vio enseguida y 
casi duplicó las métricas del estado del arte anterior.

\section{Tercera Generación: representaciones contrastivas y prior de difusión}

Las arquitecturas más nuevas dejan de lado el mapeo lineal y le dedican una
red entera al tema de extraer significado.

\subsection{MindEye y alineamiento geométrico}

MindEye~\cite{scotti2023mindeye} arranca con un diagnóstico simple: la regresión lineal no
captura la geometría del córtex visual cuando se la lleva al hiperespacio de
CLIP. La solución fue cambiar el lineal por un MLP residual entrenado con
pérdida InfoNCE bidireccional. El detalle es que el contrastivo no se conforma
con minimizar el MSE, sino que obliga al vector predicho a quedar cerca del
concepto correcto y lejos de los demás, conservando la topología relacional del
espacio CLIP.

\subsection{MindEye2}

MindEye2~\cite{scotti2024mindeye2} amplía el marco con preentrenamiento multi-sujeto y un
prior de difusión adaptativo; a día de hoy es la frontera sobre el Natural Scenes Dataset 
(NSD). En paralelo, trabajos como el benchmark NSD-Imagery~\cite{kneeland2025nsdimagery} 
empujan hacia la reconstrucción de imaginería y el modelado celular.

\section{Brecha vigente: Modality Gap y adaptadores semánticos}

De toda la revisión se desprende algo que conviene subrayar: el decodificador generativo 
(DCGAN, VQGAN o Stable Diffusion) no es el cuello de botella. Lo que sigue sin resolverse, y lo 
que esta tesis ataca, es la brecha intermodal, el Modality Gap~\cite{liang2022mind}.

Esa brecha describe el desajuste geométrico entre lo que sale del fMRI y lo que CLIP espera. A 
CLIP se lo entrenó para que imágenes y texto compartieran un sub-manifold bien armado, pero 
nadie lo preparó para alojar señales corticales dentro de él. Cuando un Ridge multivariado 
predice un vector CLIP desde los vóxeles, el estimador regularizado en $L_2$ se va hacia la 
media para bajar la varianza frente al ruido del fMRI, y termina cayendo en zonas de baja 
densidad del manifold semántico. Al recibir un embedding así, sin estructura ni norma típica, 
Stable Diffusion 2.1 unCLIP hace lo que sabe: alucina tipografía y texturas genéricas. Es, 
punto por punto, lo que veo en la corrida Ridge Lineal.

De ahí que el frente activo en neurociencia computacional sea hoy diseñar adaptadores que 
cierren esa brecha. Mi tesis se planta justo ahí, con una propuesta intermedia: un adapter 
Ridge Estocástico inspirado en MindEye y MindEye2, pero sin el entrenamiento contrastivo 
completo. La apuesta es meter variabilidad calibrada sobre la salida del Ridge Lineal para 
correr las predicciones hacia zonas más densas del manifold CLIP, y hacerlo sin pagar un MLP 
residual entero entrenado con InfoNCE. El Cap.~7 compara Ridge Lineal y Ridge Estocástico en 
lo cuantitativo y lo cualitativo; los adapters contrastivos profundos quedan como el paso 
siguiente, natural, en el Cap.~9.

\afterpage{\blankpage}


